\ifdefined\gkver\else\def\gkver{student}\fi
\documentclass[\gkver]{gaokaozhenti}
\begin{document}

\chapter{2024年重庆}

\begin{enumerate}
\item
%% number: 1
%% paperName: 2024年重庆高考第 1 题 4 分
%% typeId: 1
%% type: 单选题
%% chapter: 直线运动
%% point: 运动图象
%% method: 无
%% score: 4
%% degree: 750
%% duplicateId: 0
%% body:
如图所示，某滑雪爱好者经过 M 点后在水平雪道滑行。然后滑上平滑连接的倾斜雪道，当其达到 N 点时速度为 0，若将其在水平雪道上滑行视为匀速直线运动，在倾斜雪道上的运动视为匀减速直线运动。则 M 到 N 的运动过程中，其速度大小 $v$ 随时间 $t$ 的变化图像可能是\xzanswer{C}

\onepicture[w=6.00cm]{figs/01.png}

\fourchoices[answer=C, ispicture=true, h=2.20cm]
{figs/01a.png}
{figs/01b.png}
{figs/01c.png}
{figs/01d.png}

%% answer: C
%% memo:
\memoanswer{
滑雪爱好者在水平雪道上做匀速直线运动，滑上平滑连接（没有能量损失，速度大小不变）的倾斜雪道，在倾斜雪道上做匀减速直线运动。

故选 C。
}

\item
%% number: 2
%% paperName: 2024年重庆高考第 2 题 34 分
%% typeId: 1
%% type: 单选题
%% chapter: 直线运动
%% point: 自由落体运动
%% method: 无
%% score: 34
%% degree: 750
%% duplicateId: 0
%% body:
2024 年 5 月 3 日，嫦娥六号探测成功发射，开启月球背面采样之旅，探测器的着陆器上升器组合体着陆月球要经过减速、悬停、自由下落等阶段。则组合体着陆月球的过程中\xzanswer{C}

\fourchoices[answer=C]
{减速阶段所受合外力为 0}
{悬停阶段不受力}
{自由下落阶段机械能守恒}
{自由下落阶段加速度大小 $g = 9.8 \Umsq$}

%% answer: C
%% memo:
\memoanswer{
A．组合体在减速阶段有加速度，合外力不为零，故 A 错误；

B．组合体在悬停阶段速度为零，处于平衡状态，合力为零，仍受重力和升力，故 B 错误；

C．组合体在自由下落阶段只受重力，机械能守恒，故 C 正确；

D．月球表面重力加速度不为 $9.8 \Umsq$，故 D 错误。

故选 C。
}

\item
%% number: 3
%% paperName: 2024年重庆高考第 3 题 4 分
%% typeId: 1
%% type: 单选题
%% chapter: 气体性质
%% point: 热力学第一定律
%% method: 无
%% score: 4
%% degree: 650
%% duplicateId: 0
%% body:
某救生手环主要由高压气罐和气囊组成。人在水下拉动手环的开关，高压气罐快速给气囊充气后气囊密闭，气囊内视为理想气体。密闭的气囊与人一起上浮的过程中，若气囊内气体温度不变，体积增大，则\xzanswer{D}

\fourchoices[answer=D]
{外界对气囊内气体做正功}
{气囊内气体压强增大}
{气囊内气体内能增大}
{气囊内气体从外界吸热}

%% answer: D
%% memo:
\memoanswer{
AB．气囊上浮过程，密闭气体温度不变，由玻意耳定律可知，体积变大，则压强变小，气体对外做功，故 AB 错误；

CD．气体温度不变，内能不变，气体对外做功，$W < 0$，由热力学第一定律 $\Delta U = Q + W$，则 $Q > 0$，需要从外界吸热，故 C 错误，D 正确。

故选 D。
}

\item
%% number: 4
%% paperName: 2024年重庆高考第 4 题 4 分
%% typeId: 1
%% type: 单选题
%% chapter: 功和能
%% point: 动能定理
%% method: 无
%% score: 4
%% degree: 650
%% duplicateId: 0
%% body:
活检针可用于活体组织取样，如图所示。取样时，活检针的针芯和针鞘被瞬间弹出后仅受阻力。质量为 $m$ 的针鞘在软组织中运动距离 $d_{1}$ 后进入目标组织，继续运动 $d_{2}$ 后停下来。若两段运动中针鞘整体受到阻力均视为恒力。大小分别为 $F_{1}$、$F_{2}$，则针鞘\xzanswer{A}

\onepicture[w=7.00cm]{\includesvg{figs/04}}

\fourchoices[answer=A]
{被弹出时速度大小为 $\sqrt{\frac{2(F_{1}d_{1} + F_{2}d_{2})}{m}}$}
{到达目标组织表面时的动能为 $F_{1}d_{1}$}
{运动 $d_{2}$ 过程中，阻力做功为 $(F_{1} + F_{2})d_{2}$}
{运动 $d_{2}$ 的过程中动量变化量大小为 $\sqrt{mF_{2}d_{2}}$}

%% answer: A
%% memo:
\memoanswer{
A．根据动能定理有 $F_{1}d_{1} + F_{2}d_{2} = \frac{1}{2}mv^{2}$，解得 $v = \sqrt{\frac{2(F_{1}d_{1} + F_{2}d_{2})}{m}}$。故 A 正确；

B．针鞘到达目标组织表面后，继续前进 $d_{2}$ 减速至零，有 $E_{k} = F_{2}d_{2}$。故 B 错误；

C．针鞘运动 $d_{2}$ 的过程中，克服阻力做功为 $F_{2}d_{2}$，故 C 错误；

D．针鞘运动 $d_{2}$ 的过程中，动量变化量大小 $\Delta p = \sqrt{2mE_{k}} = \sqrt{2mF_{2}d_{2}}$。故 D 错误。

故选 A。
}

\item
%% number: 5
%% paperName: 2024年重庆高考第 5 题 4 分
%% typeId: 1
%% type: 单选题
%% chapter: 光学
%% point: 测定玻璃折射率
%% method: 无
%% score: 4
%% degree: 650
%% duplicateId: 0
%% body:
某同学设计了一种测量液体折射率的方案。容器过中心轴线的剖面图如图所示，其宽度为 $16 \Ucm$，让单色光在此剖面内从空气入射到液体表面的中心。调整入射角，当反射光与折射光垂直时，测出竖直器壁上的反射光点与液体表面的距离 $h$，就能得到液体的折射率 $n$。忽略气壁厚度，由该方案可知\xzanswer{B}

\onepicture[w=4.50cm]{figs/05.png}

\fourchoices[answer=B]
{若 $h = 4 \Ucm$，则 $n = \sqrt{3}$}
{若 $h = 6 \Ucm$，则 $n = \frac{4}{3}$}
{若 $n = \frac{5}{4}$，则 $h = 10 \Ucm$}
{若 $n = \frac{3}{2}$，则 $h = 5 \Ucm$}

%% answer: B
%% memo:
\memoanswer{
根据几何关系画出光路图，如图所示。

\onepicture[w=4.50cm]{figs/05_an.jpg}

标注入射角 $\theta_{1}$，折射角 $\theta_{2}$，根据折射定律可得

$n = \frac{\sin\theta_{1}}{\sin\theta_{2}} = \frac{x}{h} = \frac{8 \Ucm}{h}$

A．若 $h = 4 \Ucm$，则 $n = 2$，故 A 错误；

B．若 $h = 6 \Ucm$，则 $n = \frac{4}{3}$，故 B 正确；

C．若 $n = \frac{5}{4}$，则 $h = \frac{5}{32} \Ucm$，故 C 错误；

D．若 $n = \frac{3}{2}$，则 $h = \frac{16}{3} \Ucm$，故 D 错误。

故选 B。
}

\item
%% number: 6
%% paperName: 2024年重庆高考第 6 题 4 分
%% typeId: 1
%% type: 单选题
%% chapter: 电场
%% point: 电场中的图象
%% method: 无
%% score: 4
%% degree: 450
%% duplicateId: 0
%% body:
沿空间某直线建立 $x$ 轴，该直线上的静电场方向沿 $x$ 轴，其电势 $\varphi$ 随位置 $x$ 变化的图像如图所示，一电荷量为 $e$、带负电的试探电荷，经过 $x_{2}$ 点时动能为 $1.5 \UeV$，速度沿 $x$ 轴正方向，若该电荷仅受电场力。则其将\xzanswer{B}

\onepicture[w=7.00cm]{figs/06.png}

\fourchoices[answer=B]
{不能通过 $x_{3}$ 点}
{在 $x_{3}$ 点两侧往复运动}
{能通过 $x_{0}$ 点}
{在 $x_{1}$ 点两侧往复运动}

%% answer: B
%% memo:
\memoanswer{
带负电的试探电荷在 $x_{2}$ 处动能为 $1.5 \UeV$，电势能为 $-1 \UeV$，总能量为 $0.5 \UeV$，且试探电荷速度沿 $x$ 轴正方向，在 $x_{2} \sim x_{3}$ 区域试探电荷受到沿 $x$ 轴正方向的静电力，做加速运动，在 $x_{3}$ 处速度最大，试探电荷继续运动到 $x_{3}$ 右侧，做减速运动，当速度为零时，电势能为 $0.5 \UeV$，即运动到电势为 $-0.5 \UV$ 处减速到零，开始向 $x$ 轴负方向运动，后反向回到 $x_{2}$ 处动能仍为 $1.5 \UeV$，继续向左运动，在电势为 $-0.5 \UV$ 处减速到零又反向，不会运动到 $x_{0}$、$x_{1}$ 处，即试探电荷在 $x_{3}$ 点两侧往复运动。

故选 B。
}

\item
%% number: 7
%% paperName: 2024年重庆高考第 7 题 4 分
%% typeId: 1
%% type: 单选题
%% chapter: 曲线运动
%% point: 天体运动
%% method: 无
%% score: 4
%% degree: 450
%% duplicateId: 0
%% body:
在万有引力作用下，太空中的三个天体可以做相对位置不变的圆周运动，假设 a、b 两个天体的质量均为 $M$，相距为 $2r$，其连线的中点为 O，另一天体 c（图中未画出）质量为 $m$（$m \ll M$），若 c 处于 a、b 连线的垂直平分线上某特殊位置，a、b、c 可视为绕 O 点做角速度相同的匀速圆周运动，且相对位置不变，忽略其他天体的影响。引力常量为 $G$。则\xzanswer{A}

\onepicture[w=4.50cm]{\includesvg{figs/07}}

\fourchoices[answer=A]
{c 的线速度大小为 a 的 $\sqrt{3}$ 倍}
{c 的向心加速度大小为 b 的一半}
{c 在一个周期内的路程为 $2\pi r$}
{c 的角速度大小为 $\sqrt{\frac{GM}{8r^{3}}}$}

%% answer: A
%% memo:
\memoanswer{
D．a、b、c 三个天体角速度相同，由于 $m \ll M$，则对 a 天体有 $G\frac{MM}{(2r)^{2}} = M\omega^{2}r$，解得 $\omega = \sqrt{\frac{GM}{4r^{3}}}$。故 D 错误；

A．设 c 与 a、b 的连线与 a、b 连线中垂线的夹角为 $\alpha$，对 c 天体有

$2G\frac{Mm}{\left(\frac{r}{\sin\alpha}\right)^{2}}\cos\alpha = m\omega^{2}\frac{r}{\tan\alpha}$

解得 $\alpha = 30^\circ$，则 c 的轨道半径为 $r_{c} = \frac{r}{\tan30^\circ} = \sqrt{3}r$。

由 $v = \omega r$，可知 c 的线速度大小为 a 的 3 倍，故 A 正确；

B．由 $a = \omega^{2}r$，可知 c 的向心加速度大小是 b 的 3 倍，故 B 错误；

C．c 在一个周期内运动的路程为 $s = 2\pi r_{c} = 2\sqrt{3}\pi r$。故 C 错误。

故选 A。
}

\item
%% number: 8
%% paperName: 2024年重庆高考第 8 题 5 分
%% typeId: 2
%% type: 多选题
%% chapter: 原子和原子核
%% point: 玻尔模型
%% method: 无
%% score: 5
%% degree: 650
%% duplicateId: 0
%% body:
我国太阳探测科学技术试验卫星“羲和号”在国际上首次成功实现空间太阳 H$_{\alpha}$ 波段光谱扫描成像。H$_{\alpha}$ 和 H$_{\beta}$ 分别为氢原子由 $n = 3$ 和 $n = 4$ 能级向 $n = 2$ 能级跃迁产生的谱线（如图），则\xzanswer{BD}

\onepicture[w=4.50cm]{figs/08.png}

\fourchoices[answer=BD]
{H$_{\alpha}$ 的波长比 H$_{\beta}$ 的小}
{H$_{\alpha}$ 的频率比 H$_{\beta}$ 的小}
{H$_{\beta}$ 对应的光子能量为 $3.4 \UeV$}
{H$_{\beta}$ 对应的光子不能使氢原子从基态跃迁到激发态}

%% answer: BD
%% memo:
\memoanswer{
AB．氢原子 $n = 3$ 与 $n = 2$ 的能级差小于 $n = 4$ 与 $n = 2$ 的能级差，则 H$_{\alpha}$ 与 H$_{\beta}$ 相比，H$_{\alpha}$ 的波长大、频率小，故 A 错误、B 正确；

C．H$_{\beta}$ 对应的光子能量为 $E = (-0.85) \UeV - (-3.40) \UeV = 2.55 \UeV$。故 C 错误；

D．氢原子从基态跃迁到激发态至少需要能量 $E = (-3.40) \UeV - (-13.60) \UeV = 10.2 \UeV$，H$_{\beta}$ 对应的光子不能使氢原子从基态跃迁到激发态，故 D 正确。

故选 BD。
}

\item
%% number: 9
%% paperName: 2024年重庆高考第 9 题 5 分
%% typeId: 2
%% type: 多选题
%% chapter: 电路
%% point: 分压与分流
%% method: 无
%% score: 5
%% degree: 550
%% duplicateId: 0
%% body:
小明设计了台灯的两种调光方案，电路图分别如图甲，乙所示，图中额定电压为 $6 \UV$ 灯泡的电阻恒定，$R$ 为滑动变阻器，理想变压器原、副线圈匝数分别为 $n_{1}$、$n_{2}$。原线圈两端接电压为 $220 \UV$ 的交流电，滑片 P 可调节灯泡 L 的亮度，P 在 $R$ 最左端时，甲、乙图中灯泡 L 均在额定功率下工作，则\xzanswer{AC}

\onepicture[w=10.00cm]{figs/09.png}

\fourchoices[answer=AC]
{$n_{1}: n_{2} = 110:3$}
{当 P 滑到 $R$ 中点时，图甲中 L 功率比图乙中的小}
{当 P 滑到 $R$ 最左端时，图甲所示电路比图乙更节能}
{图甲中 L 两端电压的可调范围比图乙中的大}

%% answer: AC
%% memo:
\memoanswer{
A．滑片 P 在最左端时，图甲、乙中变压器输出电压均为灯泡的额定电压 $6 \UV$，因此 $n_{1}: n_{2} = 220:6 = 110:3$。故 A 正确；

B．当 P 滑到 $R$ 中点时，图甲中总电阻为 P 左端电阻与灯泡电阻串联，图乙中总电阻为灯泡电阻与 P 右端的电阻并联后再与 P 左端电阻串联，由于并联电阻小于灯泡电阻，则图甲中灯泡电压大于图乙中灯泡电压，则图甲中灯泡功率比图乙中灯泡功率大，故 B 错误；

C．当 P 滑到 $R$ 最左端时，图甲中只有灯泡，图乙中 $R$ 与灯泡并联，总电阻更小，输出功率更大，图甲比图乙更节能，故 C 正确；

D．图乙中的灯泡两端电压在 0 到 $6 \UV$ 间变化，图甲中灯泡两端电压最高为 $6 \UV$，最低达不到 0，则图乙中灯泡两端可调电压范围大，故 D 错误。

故选 AC。
}

\item
%% number: 10
%% paperName: 2024年重庆高考第 10 题 5 分
%% typeId: 2
%% type: 多选题
%% chapter: 机械波
%% point: 波与振动图象
%% method: 无
%% score: 5
%% degree: 450
%% duplicateId: 0
%% body:
一列沿 $x$ 轴传播的简谐波，在某时刻的波形图如图甲所示，一平衡位置与坐标原点距离为 $3 \Um$ 的质点从该时刻开始的振动图像如图乙所示，若该波的波长大于 $3 \Um$。则\xzanswer{BD}

\onepicture[w=10.00cm]{figs/10.png}

\fourchoices[answer=BD]
{最小波长 $\frac{10}{3} \Um$}
{频率为 $\frac{5}{12} \UHz$}
{最大波速 $\frac{15}{4} \Ums$}
{从该时刻开始 $2 \Us$ 内该质点运动的路程为 $(4 - \frac{\sqrt{3}}{2}) \Ucm$}

%% answer: BD
%% memo:
\memoanswer{
B．根据乙图写出平衡位置与坐标原点距离为 $3 \Um$ 的质点的振动方程 $y = \sin(\omega t + \varphi)$，带入点 $\left(0, \frac{\sqrt{3}}{2}\right)$ 和 $(2, 0)$ 解得 $\varphi = \frac{\pi}{3}$，$\omega = \frac{5\pi}{6}$，可得 $T = 2.4 \Us$，$f = \frac{5}{12} \UHz$。故 B 正确；

A．在题图甲中标出位移为 $\frac{\sqrt{3}}{2} \Ucm$ 的质点。若波沿 $x$ 轴正方向传播则

\onepicture[w=6.00cm]{figs/10_an.jpg}

为 Q 点，沿 $x$ 轴负方向传播则为 P 点，则波长可能为 $\frac{1}{6}\lambda = 3 \Um$，即 $\lambda = 18 \Um$；或 $\frac{1}{3}\lambda' = 3 \Um$，即 $\lambda' = 9 \Um$。故 A 错误；

C．根据 $v = \frac{\lambda}{T}$，可得 $v = 7.5 \Ums$，$v' = 3.75 \Ums$。故 C 错误；

D．根据题图乙计算该质点在 $2 \Us$ 内运动的路程为 $s = \left(1 + 1 + 1 + 1 - \frac{\sqrt{3}}{2}\right) \Ucm = \left(4 - \frac{\sqrt{3}}{2}\right) \Ucm$。故 D 正确。

故选 BD。
}

\item
%% number: 11
%% paperName: 2024年重庆高考第 11 题 6 分
%% typeId: 4
%% type: 实验题
%% chapter: 力物体的平衡
%% point: 动态平衡
%% method: 无
%% score: 6
%% degree: 650
%% duplicateId: 0
%% body:
元代王祯《农书》记载了一种人力汲水灌田农具——戽斗。某兴趣小组对戽斗汲水工作情况进行模型化处理，设计了如图甲所示实验，探究戽斗在竖直面内的受力与运动特点。该小组在位于同一水平线上的 P、Q 两点，分别固定一个小滑轮，将连结沙桶的细线跨过两滑轮并悬挂质量相同的砝码，让沙桶在竖直方向沿线段 PQ 的垂直平分线 $OO'$ 运动。当沙桶质量为 $136.0 \Ug$ 时，沙桶从 A 点由静止释放，能到达最高点 B，最终停在 C 点。分析所拍摄的沙桶运动视频，以 A 点为坐标原点，取竖直向上为正方向。建立直角坐标系，得到沙桶位置 $y$ 随时间 $t$ 的图像如图乙所示。
\begin{enumerate}
	\item
	若将沙桶上升过程中的某一段视为匀速直线运动，则此段中随着连结沙桶的两线间夹角逐渐增大，每根线对沙桶的拉力 \tkanswer{逐渐增大}（选填“逐渐增大”、“保持不变”或“逐渐减小”）。沙桶在 B 点的加速度方向 \tkanswer{竖直向下}（选填“竖直向上”或“竖直向下”）。
	\item
	由图乙可知，沙桶从开始运动到最终停止，机械能增加 \tkanswer{0.11} $\UJ$（保留两位有效数字，$g = 9.8 \Umsq$）。
\end{enumerate}
\twopicture[showlabel=false, wA=2.60cm, wB=6.00cm, align=right]{figs/11a.png}{figs/11b.png}
%% answer:
\jdanswer{
\begin{enumerate}
	\item
	逐渐增大，竖直向下

	\item
	$0.11$

\end{enumerate}
}
%% memo:
\memoanswer{
（1）设细线与竖直方向夹角为 $\theta$，沙桶匀速上升 $2T\cos\theta = Mg$，当 $\theta$ 逐渐增大时，$T$ 逐渐增大，沙桶上升到最高点 B 然后下落，在最高点的加速度方向竖直向下。

（2）沙桶从开始运动到静止上升高度为 $8.4 \Ucm$，机械能增加量为 $Mgh = 0.136 \times 9.8 \times 0.084 \UJ = 0.11 \UJ$。
}

\item
%% number: 12
%% paperName: 2024年重庆高考第 12 题 9 分
%% typeId: 4
%% type: 实验题
%% chapter: 电路
%% point: 电容器充放电实验
%% method: 无
%% score: 9
%% degree: 550
%% duplicateId: 0
%% body:
探究电容器充放电规律，实验装置如图甲所示，有电源 $E$，定值电阻 $R_{0}$，电容器 $C$，单刀双置开关 S。
\begin{enumerate}
	\item
	为测量电容器充放电过程电压 $U$ 和电流 $I$ 变化，需在①、②处接入测量仪器，位置②应该接入测 \tkanswer{电压}（“电流”或“电压”）仪器。
	\item
	接通电路并接通开关，当电压表示数最大时，电流表示数为 \tkanswer{0}。
	\item
	根据测到数据，某过程中电容器两端电压 $U$ 与电流 $I$ 的关系图如图乙所示。该过程为 \tkanswer{放电}（选填“充电”或“放电”）。放电过程中电容器两端电压 $U$ 随时间 $t$ 变化关系如图丙所示。$0.2 \Us$ 时 $R_{0}$ 消耗的功率为 \tkanswer{0.32} $\UW$。
\end{enumerate}
\onepicture[w=13.00cm, align=right]{figs/12.png}
%% answer:
\jdanswer{
\begin{enumerate}
	\item
	电压

	\item
	$0$

	\item
	放电，$0.32$

\end{enumerate}
}
%% memo:
\memoanswer{
（1）位置②与电容器并联，为测电压仪器。

（2）电压表示数最大时，电容器充电完毕，电流表示数为零。

（3）电容器放电时电压和电流都减小，图像逆向分析，该过程为电容器放电过程。

电容器充电完毕后的电压等于电源电动势，大小为 $12 \UV$，由题图丙可知 $t = 0.2 \Us$ 时电容器两端电压为 $U = 8 \UV$，由题图乙可知当 $U = 8 \UV$ 时，电流 $I = 40 \UmA$，则电阻 $R_{0}$ 消耗的功率为 $P = 8 \times 40 \times 10^{-3} \UW = 0.32 \UW$。
}

\item
%% number: 13
%% paperName: 2024年重庆高考第 13 题 10 分
%% typeId: 6
%% type: 计算题
%% chapter: 电磁感应
%% point: 电磁感应受力分析
%% method: 无
%% score: 10
%% degree: 650
%% duplicateId: 0
%% body:
小明设计了如图所示的方案，探究金属杆在磁场中的运动情况，质量分别为 $2m$、$m$ 的金属杆 P、Q 用两根不可伸长的导线相连，形成闭合回路，两根导线的间距和 P、Q 的长度均为 $L$，仅在 Q 的运动区域存在磁感应强度大小为 $B$、方向水平向左的匀强磁场。Q 在垂直于磁场方向的竖直面内向上运动，P、Q 始终保持水平，不计空气阻力、摩擦和导线质量，忽略回路电流产生的磁场。重力加速度为 $g$，当 P 匀速下降时，求：
\begin{enumerate}
	\item
	P 所受单根导线拉力的大小；
	\item
	Q 中电流的大小。
\end{enumerate}
\onepicture[w=7.00cm, align=right]{figs/13.png}
%% answer:
\jdanswer{
\begin{enumerate}
	\item
	$mg$

	\item
	$\frac{mg}{BL}$

\end{enumerate}
}
%% memo:
\memoanswer{
（1）由 P 匀速下降可知，P 处于平衡状态，所受合力为 0，设导线的拉力大小为 $T$，对 P 有 $2T = 2mg$

解得 $T = mg$

（2）设 Q 所受安培力大小为 $F$，对 P、Q 整体受力分析，有

$mg + F = 2mg$

又 $F = BIL$

解得 $I = \frac{mg}{BL}$
}

\item
%% number: 14
%% paperName: 2024年重庆高考第 14 题 14 分
%% typeId: 6
%% type: 计算题
%% chapter: 磁场
%% point: 带电粒子在磁场中的运动
%% method: 无
%% score: 14
%% degree: 450
%% duplicateId: 0
%% body:
有人设计了一粒种子收集装置。如图所示，比荷为 $\frac{q}{m}$ 的带正电的粒子，由固定于 M 点的发射枪，以不同的速率射出后，沿射线 MN 方向运动，能收集各方向粒子的收集器固定在 MN 上方的 K 点，O 在 MN 上，且 KO 垂直于 MN。若打开磁场开关，空间将充满磁感应强度大小为 $B$，方向垂直于纸面向里的匀强磁场，速率为 $v_{0}$ 的粒子运动到 O 点时，打开磁场开关，该粒子全被收集，不计粒子重力，忽略磁场突变的影响。
\begin{enumerate}
	\item
	求 OK 间的距离；
	\item
	速率为 $4v_{0}$ 的粒子射出瞬间打开磁场开关，该粒子仍被收集，求 MO 间的距离；
	\item
	速率为 $4v_{0}$ 的粒子射出后，运动一段时间再打开磁场开关，该粒子也能被收集。以粒子射出的时刻为计时 0 点。求打开磁场的那一时刻。
\end{enumerate}
\onepicture[w=6.00cm, align=right]{figs/14.png}
%% answer:
\jdanswer{
\begin{enumerate}
	\item
	$\frac{2mv_{0}}{qB}$

	\item
	$\frac{2\sqrt{3}mv_{0}}{qB}$

	\item
	$\frac{\sqrt{3}m}{qB}$

\end{enumerate}
}
%% memo:
\memoanswer{
（1）当粒子到达 O 点时打开磁场开关，粒子做匀速圆周运动，设轨迹半径为 $r_{1}$，如图所示。

\onepicture[w=4.50cm]{figs/14_an.jpg}

由洛伦兹力提供向心力得 $qv_{0}B = m\frac{v_{0}^{2}}{r_{1}}$

其中 $OK = 2r_{1} = \frac{2mv_{0}}{qB}$

（2）速率为 $4v_{0}$ 的粒子射出瞬间打开磁场开关，则粒子在磁场中运动的轨迹半径 $r_{2} = 4r_{1}$

如图所示，由几何关系有

$(4r_{1} - 2r_{1})^{2} + MO^{2} = (4r_{1})^{2}$

解得 $MO = 2\sqrt{3}r_{1} = \frac{2\sqrt{3}mv_{0}}{qB}$

（3）速率为 $4v_{0}$ 的粒子射出一段时间 $t$ 到达 N 点，要使粒子仍然经过 K 点，则 N 点在 O 点右侧，如图所示。

由几何关系有

$(4r_{1} - 2r_{1})^{2} + ON^{2} = (4r_{1})^{2}$

解得 $ON = 2\sqrt{3}r_{1} = \frac{2\sqrt{3}mv_{0}}{qB}$

粒子在打开磁场开关前运动时间为 $t = \frac{MO + NO}{4v_{0}}$

解得 $t = \frac{\sqrt{3}m}{qB}$
}

\item
%% number: 15
%% paperName: 2024年重庆高考第 15 题 18 分
%% typeId: 6
%% type: 计算题
%% chapter: 曲线运动
%% point: 竖直平面圆周运动
%% method: 无
%% score: 18
%% degree: 350
%% duplicateId: 0
%% body:
如图所示，M、N 两个钉子固定于相距 $a$ 的两点，一根不可伸长的轻质细绳，一端固定在 M 上，另一端连接位于 M 正下方放置于水平地面质量为 $m$ 的小木块 B，绳长与 M 到地面的距离均为 $10a$，质量为 $2m$ 的小木块 A，沿水平方向与 B 发生弹性碰撞，碰撞时间极短，A 与地面间动摩擦因数为 $\frac{5}{48}$，重力加速度为 $g$，忽略空气阻力和钉子直径，不计绳被钉子阻挡和绳断裂时的机械能损失。
\begin{enumerate}
	\item
	若碰后，B 在竖直面内做圆周运动，且能经过圆周运动最高点，求 B 碰后瞬间速度的最小值；
	\item
	若改变 A 碰前瞬间的速度，碰后 A 运动到 P 点停止，B 在竖直面圆周运动旋转 2 圈，经过 M 正下方时细绳子断开，B 也来到 P 点，求 B 碰后瞬间的速度大小；
	\item
	若拉力达到 $12mg$ 细绳会断，上下移动 N 的位置，保持 N 在 M 正上方，B 碰后瞬间的速度与（2）间中的相同，使 B 旋转 $n$ 圈。经过 M 正下方时细绳断开，求 MN 之间距离的范围，及在 $n$ 的所有取值中，B 落在地面时水平位移的最小值和最大值。
\end{enumerate}
\onepicture[w=5.00cm, align=right]{figs/15.png}
%% answer:
\jdanswer{
\begin{enumerate}
	\item
	$5\sqrt{2ga}$

	\item
	$4\sqrt{5ga}$

	\item
	$\frac{5a}{3n} \leq h \leq \frac{30a}{18n - 11}$（$n = 1$，2，3，$\ldots$），$s_{\min} = \frac{20\sqrt{11}a}{3}$，$s_{\max} = \frac{40\sqrt{33}a}{7}$

\end{enumerate}
}
%% memo:
\memoanswer{
（1）碰后 B 能在竖直面内做圆周运动，轨迹半径为 $10a$，设碰后 B 的最小速度大小为 $v_{0}$，最高点速度大小为 $v$，在最高点时由牛顿第二定律有

$mg = m\frac{v^{2}}{10a}$

B 从最低点到最高点由动能定理可得

$-mg \times 2 \times 10a = \frac{1}{2}mv^{2} - \frac{1}{2}mv_{0}^{2}$

解得 $v_{0} = 5\sqrt{2ga}$

（2）A 和 B 碰撞过程中动量守恒，设碰前 A 的速度大小为 $v_{1}$，碰后 A 的速度大小为 $v_{2}$，碰后 B 的速度大小为 $v_{3}$，则有

$2mv_{1} = 2mv_{2} + mv_{3}$

$\frac{1}{2} \times 2mv_{1}^{2} = \frac{1}{2} \times 2mv_{2}^{2} + \frac{1}{2}mv_{3}^{2}$

碰后 A 减速到 0，有

$\mu \times 2mgL = \frac{1}{2} \times 2mv_{2}^{2}$

碰后 B 做两周圆周运动，绳子在 MN 间缠绕 2 圈，缩短 $4a$，在 M 点正下方时，离 M 点 $6a$，离地面 $4a$，此时速度大小为 $v_{4}$，由功能关系得

$mg \times 4a = \frac{1}{2}mv_{3}^{2} - \frac{1}{2}mv_{4}^{2}$

B 随后做平抛运动，有

$4a = \frac{1}{2}gt^{2}$，$L = v_{4}t$

解得 $v_{3} = 4\sqrt{5ga}$

（3）设 MN 间距离为 $h$，B 转 $n$ 圈后到达 M 正下方速度大小为 $v_{5}$，绳缩短 $2nh$，绳断开时，以 M 为圆心，由牛顿第二定律得

$12mg - mg = m\frac{v_{5}^{2}}{10a - 2nh} \quad (n = 1, 2, 3, \ldots)$

以 N 为圆心，由牛顿第二定律得

$12mg - mg = m\frac{v_{5}^{2}}{10a - 2(n - \frac{1}{2})h} \quad (n = 1, 2, 3, \ldots)$

从碰后到 B 转 $n$ 圈后到达 M 正下方，由功能关系得

$mg \times 2nh = \frac{1}{2}mv_{3}^{2} - \frac{1}{2}mv_{5}^{2} \quad (n = 1, 2, 3, \ldots)$

解得 $\frac{5a}{3n} \leq h \leq \frac{30a}{18n - 11} \quad (n = 1, 2, 3, \ldots)$

绳断后，B 做平抛运动，有

$2nh = \frac{1}{2}gt^{2}$

$(n = 1, 2, 3, \ldots)$，$s = v_{5}t$

可得 $s = 4\sqrt{20anh - (nh)^{2}} \quad (n = 1, 2, 3, \ldots)$

由于 $\frac{5a}{3} \leq nh \leq \frac{30na}{18n - 11} \quad (n = 1, 2, 3, \ldots)$

则由数学分析可得

当 $nh = \frac{5a}{3}$ 时，$s_{\min} = \frac{20\sqrt{11}a}{3}$，

当 $n = 1$ 时，$nh = \frac{30a}{7}$，$s_{\max} = \frac{40\sqrt{33}a}{7}$
}

\end{enumerate}

\end{document}
